\section{Expanded experiments on trained networks}
\label{sec:expanded-experiments}

The strongest positive observation is that twenty unmodified trained digit
networks satisfy the new numerical-stability condition at \emph{every real
input}, with interval-verified matrix inequalities and an exact rational
final bound. All twenty also fit every training example with signed output
greater than $1/2$. The condition is nevertheless not universal: a separate
well-performing trained network has an interval-verified violation even
at an actual held-out image. The old sign condition and the new numerical
condition must be kept separate when interpreting these findings.

\subsection{Protocol and what is being measured}

We trained 168 networks in five retained cohorts: 80 initial synthetic
models, 32 synthetic extensions, 32 main digit models, 16 deeper digit
models, and eight smaller-step-size controls. Configurations were fixed
before training each cohort; later cohorts were chosen after some earlier
results were visible. These are exploratory experiments, not a publicly
preregistered confirmatory study. Failed fits were retained. No Bayesian
posterior was sampled: the experiments test properties of trained
parameters, not posterior prediction probabilities or statistical sample
complexity. The new numerical condition was proposed after training, and
its evaluation is a post-training audit of the resulting networks.

The architecture is $h_0=(x,1)$,
$h_s=\tanh(h_{s-1}W_s)$, and $z=h_Jv$, with no biases beyond the input
augmentation. Synthetic widths are $32$ or $64$, dimensions $8,10,12$,
and hidden depths $4,8,16,32$. Tasks include majority, parity, random
Boolean labels, and labels from a fixed small tanh teacher. Training
uses either the whole cube or a fixed subset, with all cube inputs
subsequently audited. Image models have width $64$ and hidden depths
$4,16,32,64$. Initialization uses Gaussian first and final weights and
orthogonal internal matrices, with gains $0.7$, $1.1$, or $1.5$ in the
synthetic cohorts and $1.1$ for images.

Both optimizers run for 6000 steps with $10^{-6}$ weight decay.
The code's label \texttt{sgd} means \emph{full-batch vanilla gradient
descent}, without momentum or minibatch noise. Its learning rate is
$0.05$. Adam uses learning rate $0.003$, except in the separately
declared $0.0003$ control; rates are multiplied by $0.3$ after $80\%$
of training. Synthetic losses are half the squared error, while image
losses are binary logistic. No test-driven rescaling is used for images.
Exact saved floating-point weights, data splits, protocols, source, and
detailed records accompany the report.

For clarity, write $E_v=N\|v\|_2^2$ and
\[
 \|\eta\|_*^2=\sum_{s=1}^J\frac{\|\eta_s\|_2^2}{N}
                    +|\eta_{\rm out}|^2.
\]
There are four distinct measurements:
\begin{enumerate}
\item correctness relative to the target labels;
\item the inherited global clean-margin requirement $|z(x)|>1/2$;
\item preservation of the network's own clean sign throughout radius
      $r_{\rm sign}=10^{-2}/\sqrt{1+JE_v}$;
\item numerical output change at most $1/16$ throughout
      $r_{\rm num}=10^{-3}/\sqrt{1+JE_v}$, as in
      Section~\ref{sec:numerical-stability}.
\end{enumerate}
The last condition imposes no all-input clean margin. Its likelihood still
requires $y_i z(x_i)>1/2$ for \emph{every training example}; we audit this
separately on the raw trained networks. A constant failed fit can be stable
without being eligible under this likelihood.

Positive pointwise tests use the finite-radius clean-transport certificate,
including nonlinear Taylor remainder bounds; they are not just tests of
the clean Jacobian. Unless explicitly marked interval verified, the
sufficient inequalities are evaluated in float64. Failure to certify is
reported as unresolved. Negative witnesses use explicit perturbations
strictly inside the normalized noise ball. Selected witnesses and the
twenty global positive results receive the stronger checks described below.

\subsection{Full-cube synthetic audits}

Of 112 synthetic runs, 105 fit every training label, and all 105 have
correct signed training outputs greater than $1/2$. Under the old global
sign-support condition, 52 pass a sufficient certificate on their entire
cube, 59 fail the clean-margin requirement at an explicit cube point,
and one remains unresolved. Seventy networks certify their own sign on
the whole cube, independently of the clean-margin requirement; 35 have
explicit, strictly budget-checked sign-flipping perturbations.

\begin{center}
\begin{tabular}{lrrrrr}
\toprule
Cohort & Runs & Fits & Old support & All-cube sign & Sign attacks\\
 & & & in/out/? & certified & found\\
\midrule
Full cube, $d=8$ & 24 & 23 & $23/1/0$ & 24 & 0\\
Half cube, $d=8$ & 24 & 23 & $12/12/0$ & 17 & 6\\
Initialization gains, $d=8$ & 16 & 16 & $7/9/0$ & 10 & 5\\
Subset, $d=10$ & 16 & 15 & $6/10/0$ & 10 & 5\\
Subset, $d=12$ & 8 & 8 & $2/6/0$ & 5 & 2\\
Small teacher, $d=10$ & 16 & 13 & $0/15/1$ & 0 & 14\\
Depth $32$, $d=10$ & 8 & 7 & $2/6/0$ & 4 & 3\\
\bottomrule
\end{tabular}
\end{center}
Here ``out'' means a witnessed clean-margin failure; ``?'' is not a
nonrobustness claim. Counts include failed fits.

The apparent fragility is concentrated near the decision boundary.
Across all 105 successful fits, \emph{every} cube point whose clean
absolute output exceeds $1/2$ passes the finite-radius sign certificate:
80,603 input--model pairs in total, with a nonempty qualifying set for
every model. Seven networks classify their entire held-out set correctly
but fail the global clean margin; all seven certify sign stability on
their whole cube. Two depth-16, dimension-12 full-batch-GD models achieve
held-out accuracies $99.55\%$ and $99.44\%$ yet admit sign attacks at rare
inputs with clean margins about $0.0072$ and $0.0021$.

Rescaling the readout using the global minimum margin makes 108 of the
112 networks pass the old support certificate, leaving four unresolved.
This is a changed-network diagnostic, not a raw training outcome:
rescaling raises the margin and simultaneously shrinks the required noise
radius, and may incur a large Gaussian energy cost. On a finite cube,
eventual membership under sufficiently large rescaling is possible for
every network with no zero clean outputs, so this statistic alone provides
little evidence of a natural training bias.

At the new numerical radius, 108 of 112 raw networks pass on the entire
cube, including \emph{all 105} training-margin-eligible fits. Of 90,112
input--model pairs, 90,108 pass; the four unresolved points belong to
failed training runs. These positive counts use floating-point evaluation
of sufficient bounds, not interval arithmetic.

\subsection{Actual images, depth, and failed optimization}

We use the locally bundled \texttt{sklearn} handwritten-digit dataset:
1797 images of $8\times8$ pixels, with tasks ``3 versus 8'' and
``even versus odd''. A $70/30$ split is stratified by the original digit
labels. Boolean features threshold pixels at $8$; continuous features
map a pixel value $p$ to $p/8-1$. These are distinct experiments. The
continuous feature results extend the numerical audit, not the Boolean
input premise of the statistical theorem. The main cohort crosses both
tasks, both feature maps, depths $4,16$, both optimizers, and two seeds.
All 32 models satisfy the training-margin likelihood.

\begin{center}
\begin{tabular}{llrrr}
\toprule
Task & Features & Test accuracy & Own-sign certificate & Margin and sign\\
\midrule
3 versus 8 & Boolean & $95.37$--$97.22\%$ & $100\%$ & $98.15$--$100\%$\\
3 versus 8 & Continuous & $97.22$--$100\%$ & $100\%$ & $99.07$--$100\%$\\
Even/odd & Boolean & $93.33$--$95.00\%$ & $99.63$--$100\%$ & $98.89$--$100\%$\\
Even/odd & Continuous & $97.04$--$98.15\%$ & $99.81$--$100\%$ & $98.89$--$100\%$\\
\bottomrule
\end{tabular}
\end{center}
Each row contains eight models; ranges are across those models. The last
two columns concern \emph{held-out image inputs} and the old radius,
not all possible inputs. Every held-out point with $|z|>1/2$ passes
the sign certificate in all 32 models. Thresholding can give identical
feature vectors in different images: at most 23 held-out vectors per
model occur in the training set. Counts and accuracy restricted to novel
feature vectors are recorded separately; these are descriptive robustness
experiments, not a benchmark generalization claim.

The separate depth-$32/64$ cohort uses only the 3-versus-8 task. All eight
full-batch-GD models fit the training likelihood and achieve
$96.30$--$100\%$ test accuracy. All eight standard-rate Adam runs fail
to fit: depth-32 accuracies range from $50.93\%$ to $75.93\%$, and all
four depth-64 runs collapse to $50.93\%$. The separately declared
smaller-rate Adam cohort is retained alongside these failures: seven of
eight fit every training margin and achieve $96.30$--$100\%$ test accuracy;
one depth-64 Boolean-feature run still collapses to $50.93\%$.
Stable constant predictions in failed runs do not support a claim about
successful learning. No failed optimization run was replaced by a
successful rerun in the reported cohort.

\subsection{Off-data sign failures and global numerical stability}

For each of the sixteen Boolean-feature models in the main digit cohort,
we sampled 10,000 uniform 64-bit inputs. All sixteen have explicit
clean-margin violations and sign-flipping perturbations at such inputs;
all sixteen attacks still work after deleting scalar output noise.
These inputs generally do not resemble digit images. A separate attack
on 32 sampled points just above clean margin $1/2$ per model found no
flips; a failed attack is not a certificate.

Three hidden-only sign witnesses are checked using exact rational noise
budgets and 60-digit interval forward evaluation of the saved binary
floating-point weights and perturbations. Their outputs change as follows:
\[
 \begin{array}{c|rr}
 \text{Model ID}&\text{clean output}&\text{perturbed output}\\\hline
 203&+0.00164981&-0.0274141\\
 204&+0.0287959&-0.364879\\
 205&-0.0192534&+0.300013
 \end{array}
\]
All three already fail the inherited $1/2$ margin. They therefore do not
show that added internal robustness excludes a parameterization which
otherwise meets the original global-margin support restriction.

The new numerical condition gives a substantially better positive result.
All 10,368 held-out input--model pairs in the main digit cohort pass its
sufficient inequality. Twenty of these 32 raw trained models---all sixteen
depth-4 models and four depth-16 models---have a \emph{global} positive
certificate, covering every real input, not only images or the cube.
Ten of these twenty use Boolean features and ten use continuous features.
These twenty results are additionally interval verified. For each
hidden-to-hidden matrix, an exact dyadic-integer Gram matrix and a
90-digit interval $LDL^\top$ decomposition prove
$W^\top W\prec g^2I$ for an explicitly stored rational $g^2$.
Cauchy--Schwarz and global Lipschitzness of tanh bound the simultaneous
noise gain. The final comparison to $(1/16)^2$ is an exact rational
inequality. The resulting output-change upper bounds range from
$0.00453$ to $0.04950$, all below $1/16$. This is a numerical proof
trusting the interval-arithmetic implementation, not a Lean formalization.

The radius calibration matters. Reusing these same verified operator bounds
and rescaling the final inequality with exact rational arithmetic, sixteen,
twelve, and four of these twenty models also certify output change at most
$1/16$ when the radius coefficient $0.001$ is increased to $0.002$, $0.005$,
and $0.01$, respectively. The last coefficient equals the old sign-radius
coefficient, although numerical stability and sign stability remain
different properties. A failed sufficient inequality at a larger radius
does not establish a violation. The rescaling code and records are supplied
separately from the frozen experiment archive.

\begin{center}
\small\begin{tabular}{lrrrr}
\toprule
Cohort & Models & Training-margin & Numerical certificate & Full tested set\\
 & & eligible & input--model pairs & certified models\\
\midrule
Synthetic, full cubes & 112 & 105 & $90{,}108/90{,}112$ & 108\\
Main digits, held out & 32 & 32 & $10{,}368/10{,}368$ & 32\\
Deep digits, held out & 16 & 8 & $1{,}728/1{,}728$ & 16\\
Smaller-rate digits, held out & 8 & 7 & $863/864$ & 7\\
\bottomrule
\end{tabular}
\end{center}
The deep numerical certificates in this table concern held-out images,
not global input support. Three failed synthetic fits and eight failed
standard-rate deep image fits pass their entire tested set; this illustrates
why the likelihood eligibility column is essential. The one unresolved
smaller-rate image point has the negative witness described next.
Pooled input--model counts reuse inputs across models and are not counts
of independent test examples for statistical inference.

Model 203 illustrates the distinction particularly clearly. It has the
old hidden-only sign flip displayed above, but the new global certificate
bounds output changes by $0.004744$ throughout the smaller numerical
radius, for every real input. Numerical stability intentionally allows
small-margin decisions to change sign; it need not rule out otherwise
useful trained networks because of a decision boundary away from the data.

There are also genuine counterexamples to automatic numerical stability.
The smaller-rate Adam model 400, with depth 32 and Boolean pixels, fits
every training margin and achieves $98.15\%$ test accuracy. At held-out
image index 25, an exactly budget-checked perturbation changes its output
from about $+0.188308$ to $-0.006960$, an absolute change of $0.195268$.
A 120-digit interval calculation proves that this exceeds $1/16$.
The true label is $-1$, so this particular perturbation \emph{corrects}
an existing classification error; it is a counterexample to numerical
stability, not an induced label error. Other verified violations occur
on uniform-cube inputs. One changes $5.296417$ to $5.218443$, which violates
raw numerical stability while leaving both clipped scores equal to one.
This supports considering the bounded-score condition of
Section~\ref{sec:clipped-stability}; it does not establish that every
trained model satisfies that weaker condition. In fact, clipping leaves
model 400's small-output held-out violation unchanged, so that trained
model also violates the bounded-score condition at the same radius.

\subsection{Other diagnostics and limits of the evidence}

The product of layer Frobenius norms can be enormous even when the new
stability test succeeds. For raw main digit models, the median
$\log_{10}(\sqrt d\prod_s\|W_s\|_F)$, including the readout as its last
factor, is $4.80$ at depth 4 and $15.74$ at depth 16. In the standard
deeper cohort it is $31.34$ at depth 32 and $58.96$ at depth 64. Thus a
small product-Frobenius condition can be much more restrictive on these
trained parameterizations than finite-noise output stability. These finite
depths do not establish an asymptotic exponential law, and the experiment
does not exclude a lower-product implementation of the same function.

A separate, fixed activation-rounding diagnostic rounds every hidden
activation and the final output to $8$, $12$, or $16$ binary fractional
bits, without rounding weights. Across the main digit models, 8-bit
rounding changes two held-out signs; 12- and 16-bit rounding change none.
Only 16 bits uniformly place the worst-case rounding budget within the
old radius for all 32 models. Success of one deterministic rounding rule
is weaker than robustness against every allowed perturbation.

These are modest tanh multilayer perceptrons, not evidence about large
contemporary ResNets or Transformers. The strongest positive global
certificates occur at depths 4 and 16; favorable on-image behavior at
depths 32 and 64 does not prove all-input membership. None of these finite
experiments proves a polynomial-depth scaling law or a universal property
of gradient descent. They support the naturalness of numerical stability
for some successful trained networks, expose its exceptions, and show why
training correctness, confident-point stability, global margin, and global
numerical stability should be reported separately.
